% GENERATED FILE. Edit canonical hl-activity contracts, then run npm run generate:books.
% canonical-activities: 8

\chapter*{Review Questions}
\addcontentsline{toc}{chapter}{Review Questions}
\markboth{Review Questions}{Review Questions}

Try these without looking ahead. Every prompt comes from the same canonical lesson data as its chapter. Answers begin in the next section.

\section*{Chapter 1}

\noindent\begin{minipage}[t]{\linewidth}
\raggedright
\hypertarget{review-JA-C01-hai-beats}{\textbf{1.1}} \emph{\ja{はい} — yes, in two equal beats}\par
\small Write the Japanese word for 'yes' in hiragana.
\end{minipage}\par\medskip

\noindent\begin{minipage}[t]{\linewidth}
\raggedright
\hypertarget{review-JA-C01-iie-length}{\textbf{1.2}} \emph{\ja{いいえ} — no, and why the extra beat matters}\par
\small How many beats (morae) does \ja{いいえ} hold?
\end{minipage}\par\medskip

\noindent\begin{minipage}[t]{\linewidth}
\raggedright
\hypertarget{review-JA-C01-konnichiwa-particle}{\textbf{1.3}} \emph{\ja{こんにちは} — hello, and a sign that lies about its sound}\par
\small The greeting ends in the sign \ja{は}. How is it pronounced there?
\end{minipage}\par\medskip

\noindent\begin{minipage}[t]{\linewidth}
\raggedright
\hypertarget{review-JA-C01-nihongo-readings}{\textbf{1.4}} \emph{\ja{日本語} — three characters, and no single sound each}\par
\small How is the character \ja{日} read inside \ja{日本語}?
\end{minipage}\par\medskip

\noindent\begin{minipage}[t]{\linewidth}
\raggedright
\hypertarget{review-JA-C01-koohii-bar}{\textbf{1.5}} \emph{\ja{コーヒー} — the third system, and a word that really is a cousin}\par
\small What does the bar \ja{ー} do in \ja{コーヒー}?
\end{minipage}\par\medskip

\noindent\begin{minipage}[t]{\linewidth}
\raggedright
\hypertarget{review-JA-C01-arigatou-dakuten}{\textbf{1.6}} \emph{\ja{ありがとう} — thanks, from “this hardly ever happens”}\par
\small Which hiragana sign becomes \ja{が} when the dakuten is added?
\end{minipage}\par\medskip

\noindent\begin{minipage}[t]{\linewidth}
\raggedright
\hypertarget{review-JA-C01-gozaimasu-level}{\textbf{1.7}} \emph{\ja{ありがとうございます} — politeness you cannot leave out}\par
\small You are thanking a shop assistant you have never met. Which form do you use?
\end{minipage}\par\medskip

\noindent\begin{minipage}[t]{\linewidth}
\raggedright
\hypertarget{review-JA-C01-practice-final-line}{\textbf{1.8}} \emph{Practice — one daytime exchange, three scripts}\par
\small Someone has just thanked you politely. Give the modest one-word reply this chapter taught.
\end{minipage}\par\medskip

\chapter*{Answer Key}
\addcontentsline{toc}{chapter}{Answer Key}
\markboth{Answer Key}{Answer Key}

The first response is the canonical display answer. When a question accepts other authored forms, they appear underneath.

\section*{Chapter 1}

\noindent\begin{minipage}[t]{\linewidth}
\raggedright
\hyperlink{review-JA-C01-hai-beats}{\textbf{1.1}} \emph{\ja{はい} — yes, in two equal beats}\par
\small \textbf{Answer:} \ja{はい}\par
\footnotesize \textbf{Also accepted:} hai; ha-i
\end{minipage}\par\medskip

\noindent\begin{minipage}[t]{\linewidth}
\raggedright
\hyperlink{review-JA-C01-iie-length}{\textbf{1.2}} \emph{\ja{いいえ} — no, and why the extra beat matters}\par
\small \textbf{Answer:} three\par
\footnotesize \textbf{Also accepted:} 3; three beats; three morae
\end{minipage}\par\medskip

\noindent\begin{minipage}[t]{\linewidth}
\raggedright
\hyperlink{review-JA-C01-konnichiwa-particle}{\textbf{1.3}} \emph{\ja{こんにちは} — hello, and a sign that lies about its sound}\par
\small \textbf{Answer:} wa\par
\footnotesize \textbf{Also accepted:} as wa; \ja{わ}; the topic particle wa
\end{minipage}\par\medskip

\noindent\begin{minipage}[t]{\linewidth}
\raggedright
\hyperlink{review-JA-C01-nihongo-readings}{\textbf{1.4}} \emph{\ja{日本語} — three characters, and no single sound each}\par
\small \textbf{Answer:} ni\par
\footnotesize \textbf{Also accepted:} as ni; \ja{に}; read ni
\end{minipage}\par\medskip

\noindent\begin{minipage}[t]{\linewidth}
\raggedright
\hyperlink{review-JA-C01-koohii-bar}{\textbf{1.5}} \emph{\ja{コーヒー} — the third system, and a word that really is a cousin}\par
\small \textbf{Answer:} it holds the vowel for one more beat\par
\footnotesize \textbf{Also accepted:} lengthens the vowel; long vowel; it lengthens the previous vowel; adds a beat to the vowel
\end{minipage}\par\medskip

\noindent\begin{minipage}[t]{\linewidth}
\raggedright
\hyperlink{review-JA-C01-arigatou-dakuten}{\textbf{1.6}} \emph{\ja{ありがとう} — thanks, from “this hardly ever happens”}\par
\small \textbf{Answer:} \ja{か}\par
\footnotesize \textbf{Also accepted:} ka; \ja{か} ka
\end{minipage}\par\medskip

\noindent\begin{minipage}[t]{\linewidth}
\raggedright
\hyperlink{review-JA-C01-gozaimasu-level}{\textbf{1.7}} \emph{\ja{ありがとうございます} — politeness you cannot leave out}\par
\small \textbf{Answer:} \ja{ありがとうございます}\par
\footnotesize \textbf{Also accepted:} arigato gozaimasu; arigatou gozaimasu; arigatō gozaimasu; the polite one
\end{minipage}\par\medskip

\noindent\begin{minipage}[t]{\linewidth}
\raggedright
\hyperlink{review-JA-C01-practice-final-line}{\textbf{1.8}} \emph{Practice — one daytime exchange, three scripts}\par
\small \textbf{Answer:} \ja{いいえ}\par
\footnotesize \textbf{Also accepted:} iie; no; not at all
\end{minipage}\par\medskip
